\documentclass[noauthor,noinstructornotes]{ximera}

\title{Change of Base}
\author{Math 1193 Design Team}

\begin{document}
\begin{abstract}
    Walking through the change of base formula for logarithms.
\end{abstract}
\maketitle

When I was learning logarithms, there was no way to put \(\log_2(5)\)
into my calculator, since there was only a \(\log\) button and a
\(\ln\) button, without the option to directly input a base other
than 10 or \(e\). In this problem, we'll learn how my generation
dealt with this problem. 

Make sure to write down your work at each step.

\begin{enumerate}
    \item Write the equation \(\log_2(5) = x\) in exponential form.
    \item Now take the \textbf{natural log} of both sides of your
    exponential equation.
    \item Next, we'll use a property of logarithms we'll learn soon:
    \(\ln(2^x) = x \ln(2)\). Use this to rewrite one side of your equation.
    \item Finally, solve for \(x\).
    \item How did I put \(\log_2(5)\) into my calculator? What about
    \(\log_7(19)\)?
\end{enumerate}

\end{document}
